\lesson{The Lean development}\label{sec:lean}
\focus{The companion instantiates the formal syntax and checks the worked calculations.}

The companion \texttt{paper/TutorialExamples.lean} formalizes the core
example using the development's \FC{} definitions. It checks the interface
coercions, store typing, views, member normalization, composition,
and function coercion.
It also constructs DOT-MNF derivations: DOT in \emph{monadic
normal form}, where intermediate computations are named by lets.
These examples check the focused translations from Section~\ref{sec:evidence}
and the complete program from Section~\ref{sec:closed-program}, including
its literal, recursive opening, lambda, and four lets.
From the repository root, run:
\begin{lstlisting}
lake env lean lean/Coercions/paper/TutorialExamples.lean
\end{lstlisting}
The imported libraries must already have been built. In a fresh checkout,
run \texttt{lake build} first.
The equalities are theorems checked by Lean's proof-checking kernel.

\subsection*{From recipe notation to constructors}
The declaration \texttt{def name : Type := expression} defines a named
expression of the stated Lean type.
Here \texttt{LeCo s} is the syntax of inclusion evidence whose variables
belong to scope \texttt{s}. This scope parameter is distinct from the
mathematical self name $s$ used in our telescopes.
The Lean type describes which expressions can be built; the evidence
typing judgment separately checks their inclusion endpoints.

Here is the companion's definition of the exposing coercion:
\begin{lstlisting}
def expose : LeCo s :=
  .obj litTel
    (.le
      (.le .nil .none (.eqSym 0) .none)
      .none (.eq 0) .none)
\end{lstlisting}
The leading dot lets Lean infer the constructor's namespace.
The \texttt{.obj} constructor combines the input telescope
\texttt{litTel} with a morphism. Reading from the inside out,
\texttt{.nil} is the empty morphism and each \texttt{.le} appends an
inclusion recipe. The inner \texttt{.le} supplies the lower recipe,
reading equality slot 0 in reverse with both sides absent.
The outer \texttt{.le} supplies the upper recipe, reading the same slot
forwards.

This is the two-row recipe table from Section~\ref{sec:recipes},
represented as an inductive list. Recovering the list order makes the
individual recipes easy to identify in the constructor expression.

Type selections are represented by scope-indexed variables and labels.
The declaration \texttt{A : Label := .typ 0} chooses a type-member label.
Within the self binder, \texttt{.here} refers to that self, so
\texttt{.sel .here A} represents $s.A$.

\subsection*{What exactly is checked?}
The companion first constructs a store containing a literal defining
$A=W$ and proves its typing. In that store, it establishes the view
equations
\[
 \view{a}=[\ev{le}\,\form{id},\ev{le}\,\form{id}],
 \qquad
 \view{b}=[\ev{le}\,\form{bot},\ev{le}\,\form{top}].
\]
The normalizer takes an explicit fuel argument; a successful result is
represented by \texttt{some} applied to the computed form.
Each example supplies a concrete fuel bound.
\texttt{CanonicalForms.lean} proves that sufficient fuel exists for
every well-typed input closed over a typed store.

The \texttt{SourceTranslation} namespace applies the translation function
to explicit source derivations. Its \texttt{nestedTranslates} and
\texttt{composedTranslates} theorems establish the nested and composed
casts; further theorems prove target typing, distinct target syntax, and
equal erasure. Those focused derivations use the parameter context of the
translation figures. The \texttt{FullSource} namespace then constructs
the closed source derivation. Its \texttt{programTranslates} theorem
identifies the complete target constructor expression,
\texttt{programTyped} derives its typing, and \texttt{programErases}
checks its erasure. \texttt{openedRoot} and \texttt{openedView} check
the two properties of $R$: preservation of the root and preservation of
the view with two additional units of fuel. The store examples check
the underlying view and normalization calculations. We do not claim a
separately mechanized multistep execution theorem for the whole tutorial
program.

The generated capture certificate is a paper calculation of the
translation clauses in \texttt{Terms.lean} and \texttt{Evidence.lean}.
It retains the generated constructor syntax. It has not been added as a
separate companion theorem.
The operational capture examples instead use hand-written target evidence.
The $h\ \ev{unit}$ trace specializes the application and prediction rules
on paper.

\needspace{5\baselineskip}
The factory follows the mechanized C1/C6 examples in
\texttt{Captures/FCdot/Examples.lean}.
It uses names for their encoded variables. Its trace groups administrative transitions; a row can
represent several machine steps.

\subsection*{Code navigation}
\begin{center}
\small
\renewcommand{\arraystretch}{1.3}
\begin{tabularx}{\linewidth}{@{}>{\raggedright\arraybackslash}p{0.35\linewidth}>{\raggedright\arraybackslash}X@{}}
\toprule
Question & Starting point, under \texttt{lean/Coercions/}\\
\midrule
What are slots and recipes? &
\leanfile{FCdot/Syntax.lean}{FCdot/Syntax.lean}; \texttt{Hole}, \texttt{Morphism}\\
Why are sides self-free? &
\leanfile{FCdot/Typing.lean}{FCdot/Typing.lean}; morphism typing\\
How are recipes combined? &
\leanfile{FCdot/Normalizer.lean}{FCdot/Normalizer.lean}; \texttt{Form.combine}, \texttt{LocalEntry.at}\\
Why do the calculations succeed? &
\leanfile{FCdot/CanonicalForms.lean}{FCdot/CanonicalForms.lean}; \texttt{le\_canon}, \texttt{atom\_canon}\\
How is a cast function called? &
\leanfile{FCdot/Machine.lean}{FCdot/Machine.lean}; application steps\\
Does the evidence checker agree? &
\leanfile{FCdot/CheckerCompleteness.lean}{FCdot/CheckerCompleteness.lean}\\
How is source safety obtained? &
\leanfile{DotToFCdot/Safety.lean}{DotToFCdot/Safety.lean}; \texttt{dot\_safety}\\
How is closing evidence generated? &
\leanfile{Captures/DotToFCdot/Terms.lean}{Captures/DotToFCdot/Terms.lean};
\texttt{translate}, \texttt{translateUses}\\
What do captures mean operationally? &
\leanfile{Captures/FCdot/Prediction.lean}{Captures/FCdot/Prediction.lean};
\texttt{app\_uses}, \texttt{capture\_prediction}, \texttt{effect\_safety}\\
Where is the factory example? &
\leanfile{Captures/FCdot/Examples.lean}{Captures/FCdot/Examples.lean};
\texttt{C1\_typed}, \texttt{C6\_run}, \texttt{C6\_safe}\\
\bottomrule
\end{tabularx}
\end{center}
\normalsize

The reference report \emph{Explicit Evidence for Dependent Object
Types} gives the full syntax and core metatheory, including the target
typing rules and canonical-forms theorem used here. The capture report develops use
sets, scopes, and classifiers. Together, the reports and the table above
connect each worked construction to its general rule and Lean declaration.

\exercise{lean}{In the displayed definition of \texttt{expose}, replace
the first \texttt{eqSym 0} by \texttt{eq 0}. What two facts would the
modified recipes produce? Why would that no longer check as an
inclusion from $\ty{Lit}$ to $\ty{Exact}$?}

\lesson{Connecting the model to a compiler}\label{sec:compiler}
\focus{Explicit evidence separates source type reconstruction from validation of its results.}

A source type checker reconstructs a typing derivation by choosing rule
instances and intermediate types and inferring missing annotations.
The target checker validates the resulting evidence.
For composition, the evidence supplies both components; checking
computes their endpoints and verifies that the intermediate types agree.

The target checker is proved sound and complete for its evidence
judgment: it accepts exactly the evidence admitted by that judgment.
This completeness result concerns evidence checking. Completeness of
source inference depends on the source frontend's search procedure.
The repository's frontends produce source derivations and compose with
translation and checking; their search behavior is a separate property.

This boundary supports two experiments with a real compiler.
First, for a specified source fragment, compare the compiler's decisions
with those of the modeled frontend on generated programs and regression
tests. Retain successful derivations and translations. Disagreements
can then be classified as differences in language rules, search
limitations, encoding errors, or compiler defects.

Second, instrument the compiler to record how it discharges typing
obligations, then translate those derivations into independently checked
evidence.
The resulting guarantee is relative to the modeled language and the
translation's stated coverage. Establishing end-to-end compiler
correctness would additionally require proofs that the compiler-to-model
translation and generated code preserve the source semantics.

\subsection*{A nominal source calculus}
A nominal source calculus would add class identity, declared
inheritance, and member lookup. Its elaboration would justify these
declarations in the structural target or in an extension with explicit
nominal evidence.

An inheritance clause generates obligations for inherited members;
a type alias generates an equality obligation; a method override
generates argument and result adaptation obligations.
These extend constructions already encountered in the tutorial.
Nominal identity, initialization, dispatch, and recursive class
structure introduce further semantic requirements.

Such an elaboration must preserve typing and operational behavior,
while retaining the evidence normalization and substitution properties
used by the safety proof.
The existing path, capture, and recursive-subtyping developments provide
components for this investigation. Their integration into a realistic
Scala model remains an open development task.

\remember{The development provides an evidence language, verified
translations for specified source calculi, and safety proofs.
A compiler connection must identify its source fragment, translate its
typing obligations, and state the guarantee obtained from checking them.}
